\chapter{HASIL DAN PEMBAHASAN}

Bab ini menyajikan hasil pengolahan data lapangan mengenai strategi pemasaran Rumah Makan Nasi Gerilya pada platform GrabFood. Data yang dianalisis meliputi observasi GrabFood Nasi Gerilya dan kompetitor, wawancara pegawai Nasi Gerilya, wawancara intern \textit{business consultant} Nasi Gerilya, wawancara kompetitor, angket pelanggan, dokumentasi kinerja, serta data pra-survei sebagai konteks awal. Informasi dari setiap sumber dipadukan untuk membaca kondisi strategi pemasaran, faktor internal dan eksternal, serta alternatif strategi yang relevan bagi peningkatan penjualan. Temuan diberi kode agar setiap faktor SWOT, IFAS, EFAS, dan strategi yang dirumuskan dapat ditelusuri kembali ke sumber data.

\section{Gambaran Umum Objek Penelitian}

\subsection{Profil Rumah Makan Nasi Gerilya}
Rumah Makan Nasi Gerilya merupakan usaha nasi Padang di Kota Medan yang memanfaatkan GrabFood sebagai salah satu kanal pemasaran dan penjualan utama. Berdasarkan data pra-survei, Nasi Gerilya mulai beroperasi pada September 2024 dan menggunakan beberapa kanal online, tetapi GrabFood menjadi kanal yang paling menonjol karena diperkirakan menyumbang lebih dari 50\% pendapatan penjualan usaha. Produk utama yang ditawarkan adalah menu nasi Padang dengan berbagai pilihan lauk seperti rendang, ayam goreng, ayam pop, dendeng, kikil, ikan, kepala ikan, serta menu tambahan seperti sambal, kuah, dan lauk pendamping.

Dari sisi operasional, data pegawai menunjukkan bahwa Nasi Gerilya memiliki pembagian kerja antara kasir, checker, display, dapur, cook, helper, butcher, senter, dan leader. Pada sisi dapur, aktivitas dimulai dari pengecekan stok, persiapan bahan, pemanasan lauk, dan pengolahan menu. Pada sisi kasir dan checker, aktivitas mencakup penerimaan pesanan, pembayaran, pencetakan bill, penempelan salinan pesanan pada box, pengecekan kelengkapan, dan komunikasi dengan driver. Data intern \textit{business consultant} juga menunjukkan adanya upaya digitalisasi internal melalui sistem antrean, menu digital, stamp digital, dan absensi digital. Pembagian kerja dan pengembangan sistem ini menunjukkan bahwa Nasi Gerilya sudah memiliki dasar proses layanan, walaupun masih menghadapi kendala saat jam ramai.

\subsection{Kondisi GrabFood sebagai Kanal Penjualan}
GrabFood berfungsi sebagai etalase digital Nasi Gerilya. Pada platform tersebut, pelanggan melihat nama toko, rating, jumlah penilaian, foto menu, harga, promo, jarak, ulasan, dan estimasi pengantaran sebelum membeli. Observasi tanggal 13 Juni 2026 menunjukkan bahwa Nasi Gerilya memiliki rating 4,7 dari 4.584 penilaian, 19 promo/voucher aktif, beberapa menu berlabel terlaris, serta ulasan yang menonjolkan rasa, porsi besar, bahan segar, kemasan, variasi menu, dan harga yang dianggap terjangkau. Pada saat yang sama, ulasan negatif/netral dan data pelanggan juga menunjukkan masalah yang perlu diperhatikan, seperti item kurang lengkap, add-on atau kuah tidak diberikan, menu habis, persepsi harga mahal, dan konsistensi layanan saat ramai.

\subsection{Kondisi Kompetisi pada Platform GrabFood}
Ruang persaingan Nasi Gerilya di GrabFood cukup padat. Observasi kompetitor menunjukkan bahwa Nasi Gerilya bersaing dengan restoran Padang atau merchant sejenis yang memiliki kekuatan berbeda-beda. Restoran Paripurna memiliki rating 4,7 dari 1.264 penilaian dan jarak 3,0 km; Rumah Makan Dendeng Batokok memiliki rating 4,8 dari 1.479 penilaian dan jarak 7,3 km; Rumah Makan Istana Krakatau memiliki rating 4,6 dari 6.256 penilaian dan jarak 8,4 km; RM. Pondok Gurih memiliki rating 4,9 dari 56.157 penilaian dan jarak 5,0 km; sedangkan Restoran Garuda memiliki rating 4,6 dari 17.068 penilaian dan jarak 6,7 km. Posisi Nasi Gerilya cukup kuat dari sisi rating, tetapi masih menghadapi kompetitor yang memiliki jumlah penilaian lebih besar, rating lebih tinggi, jarak lebih dekat, atau reputasi lebih kuat.

\section{Deskripsi Data Penelitian}
Data penelitian utama terdiri atas wawancara pegawai Nasi Gerilya, wawancara intern \textit{business consultant}, wawancara kompetitor, angket pelanggan, observasi GrabFood, dan dokumentasi kinerja. Ringkasan data disajikan pada Tabel~\ref{tab:ringkasan-data-penelitian-utama-bab4}.

\begingroup
    \footnotesize
    \setlength{\tabcolsep}{2pt}
    \renewcommand{\arraystretch}{1.05}
    \begin{longtable}{|Z{2.4cm}|L{3.1cm}|L{2.9cm}|L{4.7cm}|}
        \caption{Ringkasan Data Penelitian Utama}
        \label{tab:ringkasan-data-penelitian-utama-bab4}\\
        \toprule
        \textbf{Jenis Data} & \textbf{Sumber} & \textbf{Cakupan Data} & \textbf{Kegunaan Analisis} \\ \hline
        \midrule
        \endfirsthead
        \multicolumn{4}{@{}l@{}}{\footnotesize\textit{Lanjutan Tabel \thetable}}\\
        \toprule
        \textbf{Jenis Data} & \textbf{Sumber} & \textbf{Cakupan Data} & \textbf{Kegunaan Analisis} \\ \hline
        \midrule
        \endhead
        \bottomrule
        \multicolumn{4}{r}{\footnotesize\textit{Bersambung ke halaman berikutnya}}\\
        \endfoot
        \bottomrule
        \multicolumn{4}{@{}L{\textwidth}@{}}{\footnotesize Sumber: Diolah peneliti berdasarkan data penelitian utama (2026).}\\
        \endlastfoot
        Wawancara pegawai & Bella sebagai pihak dapur/leader dan Ririn sebagai kasir/checker & Kode temuan TW-KR dan TW-AD & Menjelaskan alur operasional, kekuatan produk, SOP, stok, packing, driver, promo, dan kendala layanan. \\ \hline
        Wawancara intern \textit{business consultant} & Mahasiswa magang IT/Sistem Informasi melalui program APINDO UMKM Merdeka & Kode temuan W-BC-01--W-BC-21 & Menjelaskan kebutuhan digitalisasi internal, sistem antrean, menu digital, stamp digital, target pelanggan, positioning, indikator evaluasi, dan prioritas strategi operasional. \\ \hline
        Wawancara kompetitor & Pondok Krakatau dan Istana Krakatau & Kode temuan TW-KP-PK dan TW-KP-IK & Menjadi pembanding strategi harga, promo, rasa, layanan cepat, stok, packing, dan pembeda kompetitor. \\ \hline
        Angket pelanggan & Reyza, Alip, dan Ojan & Kode temuan TS-PL & Menggali alasan membeli, menu favorit, persepsi harga/promo, porsi, kemasan, keluhan, repeat order, dan pembanding kompetitor. \\ \hline
        Observasi GrabFood & Nasi Gerilya tanggal 13 Juni 2026 & Kode temuan TO-GF-01--TO-GF-08 & Menunjukkan reputasi digital, promo, menu, harga, ketersediaan, dan ulasan pelanggan. \\ \hline
        Observasi kompetitor & Lima merchant pembanding tanggal 19--20 Juni 2026 & Kode temuan TO-KP-01--TO-KP-07 & Menjelaskan posisi Nasi Gerilya terhadap rating, jumlah penilaian, jarak, harga, promo, dan tema ulasan pesaing. \\ \hline
        Dokumentasi kinerja & Data indeks penjualan, AOV, kunjungan, dan kompetitor & Indikator historis kinerja dan pembanding kompetitor & Menjadi konteks bahwa penjualan dan indeks kompetitor perlu dijelaskan melalui strategi pemasaran dan pengalaman pelanggan. \\ \hline
    \end{longtable}
\endgroup

\subsection{Konteks Pra-Survei dalam Analisis}
Selain data penelitian utama, analisis juga menggunakan data pra-survei sebagai konteks awal. Data pra-survei tidak ditempatkan sebagai pengganti data lapangan Juni--Juli 2026, tetapi sebagai gambaran pendahuluan mengenai masalah penjualan, perilaku pelanggan, dan posisi Nasi Gerilya pada kanal digital. Pemetaan pra-survei terdiri atas 128 kode, yaitu 24 kode wawancara awal pemilik, 39 kode observasi GrabFood awal, 41 kode dokumentasi kinerja, dan 24 kode foto produk/storefront. Ringkasan pemakaian data pra-survei disajikan pada Tabel~\ref{tab:konteks-prasurvei-bab4}.

\begingroup
    \scriptsize
    \setlength{\tabcolsep}{2pt}
    \renewcommand{\arraystretch}{1.02}
    \begin{longtable}{|Z{1.8cm}|L{4.2cm}|L{3.7cm}|L{3.3cm}|}
        \caption{Konteks Pra-Survei dalam Analisis Bab IV}
        \label{tab:konteks-prasurvei-bab4}\\
        \toprule
        \textbf{Kode Sumber} & \textbf{Temuan Penting} & \textbf{Fungsi dalam Analisis} & \textbf{Peran Analisis} \\ \hline
        \midrule
        \endfirsthead
        \multicolumn{4}{@{}l@{}}{\footnotesize\textit{Lanjutan Tabel \thetable}}\\
        \toprule
        \textbf{Kode Sumber} & \textbf{Temuan Penting} & \textbf{Fungsi dalam Analisis} & \textbf{Peran Analisis} \\ \hline
        \midrule
        \endhead
        \bottomrule
        \multicolumn{4}{r}{\footnotesize\textit{Bersambung ke halaman berikutnya}}\\
        \endfoot
        \bottomrule
        \multicolumn{4}{@{}L{\textwidth}@{}}{\footnotesize Sumber: Diolah peneliti berdasarkan koding pra-survei (2026).}\\
        \endlastfoot
        PS-WPM & GrabFood disebut sebagai kanal utama dengan estimasi kontribusi lebih dari 50\% penjualan; pelanggan baru naik, tetapi repeat order dan reactivated user turun; dendeng dan ayam pop putih muncul sebagai produk unggulan; promo berpengaruh tetapi harus dihitung dengan pajak, HPP, dan margin; masalah stok, antrean driver, konsistensi rasa, dan pesanan tidak sesuai request sudah muncul sejak pra-survei. & Menguatkan fokus penelitian pada GrabFood, isu loyalitas, kebutuhan hero menu, promo terukur, update stok, dan SOP peak hour. & Menjadi konteks awal untuk membaca kontribusi GrabFood, pertimbangan margin, dan pola masalah layanan. \\ \hline
        PS-GF & Observasi 4 April 2026 menunjukkan rating 4,7 dari 4.361 penilaian, variasi menu, promo, add-on, lauk tambahan, sambal kemasan, kuah/minuman, dan menu premium. & Menjadi pembanding historis terhadap observasi 13 Juni 2026, terutama untuk reputasi digital, rentang produk, promo, dan peluang upselling. & Angka terbaru yang dipakai untuk temuan utama tetap observasi 13 Juni 2026. \\ \hline
        PS-KIN & Data kinerja awal menunjukkan penjualan bulanan mencapai puncak pada Maret 2025 lalu cenderung melemah sampai November 2025; indeks Nasi Gerilya turun dari 100,0 menjadi 70,4 pada Juli--Oktober 2025 atau berubah -29,6\%; AOV dan kunjungan naik pada sebagian periode tetapi penjualan dan indeks relatif melemah. & Menjelaskan urgensi strategi peningkatan penjualan, kebutuhan menjaga konversi, repeat order, harga, promo, dan pengalaman layanan. & Menjadi indikator historis untuk membaca urgensi strategi, bukan ukuran dampak setelah rekomendasi diterapkan. \\ \hline
        PS-PROD & Foto produk dan storefront menunjukkan bukti visual variasi menu, termasuk ayam pop putih dan nasi dendeng sebagai menu unggulan. & Mendukung analisis \textit{product} dan \textit{physical evidence}, terutama penguatan foto, deskripsi, menu hero, dan tampilan digital. & Foto produk membuktikan ketersediaan aset visual, bukan membuktikan menu paling laku. \\ \hline
    \end{longtable}
\endgroup

Dengan demikian, temuan Bab IV dibaca sebagai analisis kualitatif berbasis triangulasi. Strategi yang dirumuskan diposisikan sebagai alternatif yang berpotensi mendukung peningkatan penjualan melalui perbaikan keputusan pembelian, pengalaman layanan, repeat order, dan efisiensi proses. Data pra-survei dipakai sebagai konteks historis dan penguat pola, sedangkan keputusan analisis utama disandarkan pada data lapangan Juni--Juli 2026.

\section{Hasil Wawancara, Observasi, dan Angket}

\subsection{Hasil Wawancara Pegawai Nasi Gerilya}
Wawancara pegawai menunjukkan bahwa kekuatan internal Nasi Gerilya tidak hanya berada pada rasa, tetapi juga pada adanya pembagian kerja dan prosedur dasar dalam memproses pesanan. Namun, pada saat ramai, proses tersebut masih rawan terhambat oleh human error, pesanan mendadak, keterbatasan SDM, komunikasi yang tidak selalu jelas, dan waktu packing yang membuat driver menunggu. Ringkasan koding pegawai disajikan pada Tabel~\ref{tab:koding-pegawai-bab4}.

\begingroup
    \scriptsize
    \setlength{\tabcolsep}{2pt}
    \renewcommand{\arraystretch}{1.02}
    \begin{longtable}{|Z{1.55cm}|L{5.55cm}|Z{2.75cm}|Z{2.65cm}|}
        \caption{Ringkasan Koding Wawancara Pegawai Nasi Gerilya}
        \label{tab:koding-pegawai-bab4}\\
        \toprule
        \textbf{Kode} & \textbf{Temuan} & \textbf{Tema} & \textbf{Kategori} \\ \hline
        \midrule
        \endfirsthead
        \multicolumn{4}{@{}l@{}}{\footnotesize\textit{Lanjutan Tabel \thetable}}\\
        \toprule
        \textbf{Kode} & \textbf{Temuan} & \textbf{Tema} & \textbf{Kategori} \\ \hline
        \midrule
        \endhead
        \bottomrule
        \multicolumn{4}{r}{\footnotesize\textit{Bersambung ke halaman berikutnya}}\\
        \endfoot
        \bottomrule
        \multicolumn{4}{@{}L{\textwidth}@{}}{\footnotesize Sumber: Diolah peneliti dari wawancara pegawai Nasi Gerilya (2026).}\\
        \endlastfoot
        TW-KR-01 & Dapur memiliki pembagian kerja leader, senter, butcher, cook, helper, pengecekan stok, dan persiapan bahan. & People, process & Kekuatan \\ \hline
        TW-KR-03 & Pesanan ramai pada akhir pekan dan pesanan mendadak dapat membuat proses menumpuk. & Process & Kelemahan \\ \hline
        TW-KR-04 & Keterlambatan dapat terjadi karena waktu sempit, tekanan pesanan, dan kekurangan SDM. & Process, people & Kelemahan \\ \hline
        TW-KR-05 & Dapur memakai standar resep, takaran, laporan penjualan, dan double-check rasa. & Product, process & Kekuatan \\ \hline
        TW-KR-07 & Request pelanggan hanya dapat dipenuhi dalam batas tertentu karena sistem masak massal. & Product, process & Kelemahan \\ \hline
        TW-KR-08 & Keluhan yang muncul antara lain aroma kikil, menu habis, dan risiko kualitas bahan tertentu. & Product & Kelemahan \\ \hline
        TW-AD-03 & Aplikasi online dibuka sekitar 09.45 dan lauk diperbarui sebelum pesanan masuk. & Place, process & Kekuatan \\ \hline
        TW-AD-04 & Pesanan GrabFood masuk melalui notifikasi dan print bill, lalu checker menempelkan salinan pesanan pada box. & Process & Kekuatan \\ \hline
        TW-AD-05 & Item basah/berkuah yang dipisah dari nasi rawan tertinggal saat hectic. & Process & Kelemahan \\ \hline
        TW-AD-07 & Driver diberi nomor antrean, tetapi tetap perlu estimasi tunggu yang jelas saat makanan belum selesai. & Process, people & Kekuatan/kelemahan \\ \hline
        TW-AD-09 & Harga GrabFood kadang dianggap mahal karena berbeda dari take away, tetapi promo membuat pelanggan tetap memilih GrabFood. & Price, promotion & Kelemahan/peluang \\ \hline
        TW-AD-10 & Promo dapat membuat order membeludak sehingga dapur perlu diberi informasi menu promo. & Promotion, process & Peluang/kelemahan \\ \hline
    \end{longtable}
\endgroup

Temuan pegawai memperlihatkan bahwa Nasi Gerilya sudah memiliki struktur dasar layanan online, tetapi titik kritisnya berada pada proses setelah pesanan masuk. Pesanan yang tampak sederhana pada aplikasi sebenarnya melewati beberapa tahap: update stok, penerimaan order, print bill, penempelan salinan pesanan, pembuatan makanan, pengecekan item, komunikasi dengan driver, dan penyerahan pesanan. Semakin ramai kondisi toko, semakin tinggi risiko item kurang, catatan khusus terlewat, atau driver menunggu terlalu lama. Oleh karena itu, kelemahan utama pada proses bukan ketiadaan prosedur, melainkan konsistensi pelaksanaan prosedur saat beban kerja meningkat.

\subsection{Hasil Wawancara Intern \textit{Business Consultant} Nasi Gerilya}
Wawancara intern \textit{business consultant} memperkuat temuan pegawai bahwa persoalan pemasaran Nasi Gerilya tidak hanya terletak pada promosi, tetapi juga pada kesiapan sistem internal. Informan merupakan mahasiswa magang bidang IT/Sistem Informasi yang ditempatkan melalui program APINDO UMKM Merdeka. Keterlibatannya berfokus pada pengembangan sistem antrean, absensi digital, menu digital, dan stamp digital agar operasional Nasi Gerilya lebih rapi, efektif, efisien, dan profesional. Ringkasan koding wawancara disajikan pada Tabel~\ref{tab:koding-business-consultant-bab4}.

\begingroup
    \scriptsize
    \setlength{\tabcolsep}{2pt}
    \renewcommand{\arraystretch}{1.02}
    \begin{longtable}{|Z{1.85cm}|L{5.45cm}|Z{2.85cm}|Z{2.35cm}|}
        \caption{Ringkasan Koding Wawancara Intern \textit{Business Consultant} Nasi Gerilya}
        \label{tab:koding-business-consultant-bab4}\\
        \toprule
        \textbf{Kode} & \textbf{Temuan} & \textbf{Tema} & \textbf{Kategori} \\ \hline
        \midrule
        \endfirsthead
        \multicolumn{4}{@{}l@{}}{\footnotesize\textit{Lanjutan Tabel \thetable}}\\
        \toprule
        \textbf{Kode} & \textbf{Temuan} & \textbf{Tema} & \textbf{Kategori} \\ \hline
        \midrule
        \endhead
        \bottomrule
        \multicolumn{4}{r}{\footnotesize\textit{Bersambung ke halaman berikutnya}}\\
        \endfoot
        \bottomrule
        \multicolumn{4}{@{}L{\textwidth}@{}}{\footnotesize Sumber: Diolah peneliti dari wawancara intern \textit{business consultant} Nasi Gerilya (2026).}\\
        \endlastfoot
        W-BC-01--W-BC-04 & Pendampingan berfokus pada pengembangan sistem internal berdasarkan observasi outlet, jam ramai GrabFood, pola antrean, harga jual, kebutuhan karyawan, dan diskusi dengan Nasi Gerilya. & People, process, place & Kekuatan \\ \hline
        W-BC-05 & Masalah utama adalah antrean tiga jalur: dine in, ojek online, dan take away. Saat jam makan siang, sistem antrean yang belum rapi membuat operasional mudah kacau. & Process, people & Kelemahan \\ \hline
        W-BC-05, W-BC-11, W-BC-16 & Menu digital, payment gateway, stamp digital, dan pengelolaan data pelanggan diarahkan untuk mempercepat layanan dine in/take away dan membangun loyalitas di luar platform online. & Process, physical evidence, promotion & Peluang/\allowbreak kekuatan \\ \hline
        W-BC-06--W-BC-08 & GrabFood dinilai sebagai kanal penting; target potensial adalah pekerja, mahasiswa, pelanggan sekitar, dan pelanggan jam makan siang yang mencari nasi Padang praktis, cepat, kuat rasa, dan mengenyangkan. & Place, STP, positioning & Peluang \\ \hline
        W-BC-09--W-BC-10 & Ayam pop layak dijadikan hero menu, sedangkan promo perlu dihitung agar tetap sesuai harga jual dan HPP. & Product, price, promotion & Kekuatan/\allowbreak ancaman \\ \hline
        W-BC-12--W-BC-14 & Promosi perlu mendukung sistem digital dan loyalitas; pembeda Nasi Gerilya berada pada rasa lauk, jumlah pelanggan, dan kemauan UMKM untuk berbenah melalui sistem digital. & Promotion, product, process & Kekuatan/\allowbreak peluang \\ \hline
        W-BC-15--W-BC-17 & Kendala layanan paling berat terjadi saat makan siang, hari Minggu, dan hari raya; ancaman berasal dari persaingan kuliner, biaya platform, kenaikan bahan baku, dan sensitivitas harga pelanggan. & Process, price, competition & Kelemahan/\allowbreak ancaman \\ \hline
        W-BC-18--W-BC-21 & Indikator evaluasi mencakup jumlah pesanan, omzet, waktu pelayanan, antrean tertangani, repeat order, rating, komplain, dan efektivitas sistem. Prioritas utama adalah sistem antrean, menu digital, stamp digital, dan kesiapan internal. & Control, process, loyalty & Rekomendasi strategis \\ \hline
    \end{longtable}
\endgroup

Temuan W-BC menunjukkan bahwa digitalisasi internal dapat menjadi penghubung antara strategi pemasaran dan pengalaman layanan. Sistem antrean membantu mengurangi penumpukan pesanan lintas kanal, menu digital membantu pelanggan melihat dan memesan menu secara mandiri, stamp digital membuka peluang data pelanggan dan loyalitas, sedangkan absensi digital mendukung pengelolaan internal. Artinya, promosi GrabFood akan lebih berdampak jika ditopang oleh sistem yang menjaga kecepatan, ketepatan, dan data evaluasi. Temuan ini juga memperkuat rekomendasi bahwa Nasi Gerilya perlu mengelola pelanggan di luar platform melalui database/stamp digital agar tidak sepenuhnya bergantung pada voucher dan trafik aplikasi.

\subsection{Hasil Wawancara Kompetitor}
Wawancara kompetitor memperlihatkan bahwa rumah makan pembanding juga menggunakan GrabFood, promo, dan kecepatan layanan sebagai bagian dari daya saing. Pondok Krakatau menyebut GrabFood menyumbang sekitar lebih kurang 50\% dari total pesanan, sedangkan Istana Krakatau masih didominasi offline sekitar 80\%. Kedua kompetitor sama-sama menonjolkan rasa dan pelayanan cepat. Namun, keduanya juga menghadapi masalah yang mirip, seperti driver menunggu, item tertinggal, pesanan kurang lengkap, stok belum diperbarui, dan konsistensi porsi atau rasa.

\begingroup
    \scriptsize
    \setlength{\tabcolsep}{2pt}
    \renewcommand{\arraystretch}{1.02}
    \begin{longtable}{|Z{1.75cm}|L{5.45cm}|Z{2.75cm}|Z{2.45cm}|}
        \caption{Ringkasan Koding Wawancara Kompetitor}
        \label{tab:koding-kompetitor-bab4}\\
        \toprule
        \textbf{Kode} & \textbf{Temuan} & \textbf{Tema} & \textbf{Makna SWOT} \\ \hline
        \midrule
        \endfirsthead
        \multicolumn{4}{@{}l@{}}{\footnotesize\textit{Lanjutan Tabel \thetable}}\\
        \toprule
        \textbf{Kode} & \textbf{Temuan} & \textbf{Tema} & \textbf{Makna SWOT} \\ \hline
        \midrule
        \endhead
        \bottomrule
        \multicolumn{4}{r}{\footnotesize\textit{Bersambung ke halaman berikutnya}}\\
        \endfoot
        \bottomrule
        \multicolumn{4}{@{}L{\textwidth}@{}}{\footnotesize Sumber: Diolah peneliti dari wawancara Pondok Krakatau dan Istana Krakatau (2026).}\\
        \endlastfoot
        TW-KP-PK-01 & GrabFood menyumbang sekitar lebih kurang 50\% dari total pesanan Pondok Krakatau. & Place & Ancaman/peluang \\ \hline
        TW-KP-PK-03 & Menu sering dipesan adalah ayam goreng, nasi ayam rendang, nasi ikan kakap, dan nasi ikan nila dengan harga sekitar Rp38 ribuan. & Product, price & Ancaman \\ \hline
        TW-KP-PK-04 & Pelanggan kembali karena rasa enak dan pelayanan cepat. & Product, process & Ancaman \\ \hline
        TW-KP-PK-06 & Stok habis kadang belum diperbarui di platform saat pelanggan sudah order. & Process, place & Peluang \\ \hline
        TW-KP-PK-07 & Porsi dan rasa kadang tidak konsisten karena petugas bungkus atau masak berbeda. & Product, people & Peluang \\ \hline
        TW-KP-PK-08 & Pesanan kadang tertinggal, lalu dikirim ulang melalui Grab. & Process & Peluang \\ \hline
        TW-KP-IK-01 & Penjualan Istana Krakatau masih didominasi offline sekitar 80\%. & Place & Ancaman \\ \hline
        TW-KP-IK-03 & Menu online yang disebut adalah ikan sambal dengan harga sekitar Rp37 ribuan. & Product, price & Ancaman \\ \hline
        TW-KP-IK-06 & Pesanan kurang lengkap menjadi keluhan pelanggan. & Process & Peluang \\ \hline
        TW-KP-IK-09 & Keunggulan Istana Krakatau berada pada rasa, porsi online yang lebih besar, dan empat jenis sambal terutama sambal Padang. & Product & Ancaman \\ \hline
    \end{longtable}
\endgroup

Data kompetitor menunjukkan bahwa masalah operasional online bukan hanya dialami Nasi Gerilya. Hal ini menjadi peluang karena Nasi Gerilya dapat membedakan diri melalui ketelitian packing, update stok, dan konsistensi rasa. Namun, kompetitor tetap menjadi ancaman karena mereka juga memiliki kekuatan rasa, pelayanan cepat, menu khas, promosi, reputasi offline, dan porsi yang dipersepsikan menarik.

\subsection{Hasil Angket Pelanggan}
Angket pelanggan berisi tiga responden, yaitu Reyza, Alip, dan Ojan. Ketiganya menunjukkan pola yang konsisten: pelanggan memakai GrabFood ketika membutuhkan makanan praktis, berat, dan mengenyangkan; keputusan membeli dipengaruhi oleh harga total, ongkir, promo, rating, foto, jarak, dan estimasi waktu; serta minat membeli ulang sangat dipengaruhi oleh rasa menu unggulan, porsi, promo, dan ketelitian pesanan.

\begingroup
    \scriptsize
    \setlength{\tabcolsep}{2pt}
    \renewcommand{\arraystretch}{1.02}
    \begin{longtable}{|Z{1.55cm}|L{5.55cm}|Z{2.75cm}|Z{2.55cm}|}
        \caption{Ringkasan Koding Angket Pelanggan}
        \label{tab:koding-angket-pelanggan-bab4}\\
        \toprule
        \textbf{Kode} & \textbf{Temuan} & \textbf{Tema} & \textbf{Kategori} \\ \hline
        \midrule
        \endfirsthead
        \multicolumn{4}{@{}l@{}}{\footnotesize\textit{Lanjutan Tabel \thetable}}\\
        \toprule
        \textbf{Kode} & \textbf{Temuan} & \textbf{Tema} & \textbf{Kategori} \\ \hline
        \midrule
        \endhead
        \bottomrule
        \multicolumn{4}{r}{\footnotesize\textit{Bersambung ke halaman berikutnya}}\\
        \endfoot
        \bottomrule
        \multicolumn{4}{@{}L{\textwidth}@{}}{\footnotesize Sumber: Diolah peneliti dari angket pelanggan GrabFood Nasi Gerilya (2026).}\\
        \endlastfoot
        TS-PL-01 & Pelanggan memakai GrabFood saat lelah, jam kantor padat, malam setelah tugas, atau malas keluar. & Segmentasi, place & Peluang \\ \hline
        TS-PL-02 & Pelanggan memilih Nasi Gerilya karena menu, rating, foto, jarak, dan estimasi yang masuk akal. & Targeting, physical evidence & Kekuatan/peluang \\ \hline
        TS-PL-03 & Dendeng bakar, ayam pop creamy, dan ayam goreng bumbu menjadi menu yang paling menonjol. & Product, positioning & Kekuatan \\ \hline
        TS-PL-04 & Porsi nasi dinilai besar dan mengenyangkan. & Product, price & Kekuatan \\ \hline
        TS-PL-05 & Harga dirasakan agak mahal atau di atas rata-rata, tetapi masih diterima jika ada promo, ongkir murah, atau kualitas stabil. & Price, promotion & Kelemahan/peluang \\ \hline
        TS-PL-06 & Promo, voucher, dan ongkir murah sangat memengaruhi keputusan beli. & Promotion, price & Peluang/ancaman \\ \hline
        TS-PL-07 & Foto, rating, nama menu, dan deskripsi membantu, tetapi detail isi paket, kuah, sambal, stok, dan paket hemat perlu diperjelas. & Physical evidence & Kekuatan/kelemahan \\ \hline
        TS-PL-08 & Kemasan dinilai cukup rapi dan bersih, tetapi segel, label, dan pemisahan item berkuah masih perlu ditingkatkan. & Physical evidence, process & Kekuatan/kelemahan \\ \hline
        TS-PL-09 & Salah satu pelanggan pernah mengalami item kurang, tetapi restoran mengirim ulang sehingga masalah selesai. & Process, people & Kelemahan/kekuatan \\ \hline
        TS-PL-11 & Alasan membeli ulang adalah rasa, porsi, rating, promo, dan respons keluhan; alasan tidak membeli ulang adalah harga tinggi, voucher tidak ada, pelayanan buruk, item kurang, atau rasa tidak konsisten. & Loyalty, product, price & Kekuatan/kelemahan \\ \hline
    \end{longtable}
\endgroup

Pelanggan tidak hanya menilai produk dari rasa. Dalam konteks GrabFood, pelanggan membaca nilai dari kombinasi foto, rating, ulasan, harga total, promo, ongkir, estimasi waktu, dan pengalaman saat pesanan diterima. Reyza menonjolkan dendeng bakar, porsi besar, jarak dekat, dan promo; Alip menonjolkan ayam pop, rating, respons restoran terhadap item kurang, dan kebutuhan paket kantor; sedangkan Ojan menonjolkan ayam goreng, porsi besar, voucher, dan kebutuhan paket anak kos. Pola ini menunjukkan bahwa strategi pemasaran Nasi Gerilya perlu diarahkan pada menu unggulan yang jelas, paket sesuai segmen, harga yang terbaca masuk akal, serta proses yang teliti.

\subsection{Hasil Observasi GrabFood dan Kompetitor}
Observasi GrabFood pada 13 Juni 2026 dilakukan terhadap halaman toko Nasi Gerilya - Glugur Kota. Data yang dikodekan mencakup profil toko, rating, jumlah penilaian, promo atau voucher, menu dan harga, menu yang tidak tersedia, rentang harga, serta ulasan pelanggan positif dan negatif/netral. Ringkasan temuan observasi Nasi Gerilya disajikan pada Tabel~\ref{tab:ringkasan-observasi-grabfood-20260613-bab4}.

\begingroup
    \footnotesize
    \setlength{\tabcolsep}{2pt}
    \renewcommand{\arraystretch}{1.05}
    \begin{longtable}{|Z{1.55cm}|Z{2.85cm}|L{8.75cm}|}
        \caption{Ringkasan Hasil Observasi GrabFood Nasi Gerilya 13 Juni 2026}
        \label{tab:ringkasan-observasi-grabfood-20260613-bab4}\\
        \toprule
        \textbf{Kode} & \textbf{Aspek} & \textbf{Temuan Observasi} \\ \hline
        \midrule
        \endfirsthead
        \multicolumn{3}{@{}l@{}}{\footnotesize\textit{Lanjutan Tabel \thetable}}\\
        \toprule
        \textbf{Kode} & \textbf{Aspek} & \textbf{Temuan Observasi} \\ \hline
        \midrule
        \endhead
        \bottomrule
        \multicolumn{3}{r}{\footnotesize\textit{Bersambung ke halaman berikutnya}}\\
        \endfoot
        \bottomrule
        \multicolumn{3}{@{}L{\textwidth}@{}}{\footnotesize Sumber: Observasi GrabFood Nasi Gerilya pada 13 Juni 2026; diolah peneliti (2026).}\\
        \endlastfoot
        TO-GF-01 & Identitas dan reputasi toko & Nama toko Nasi Gerilya - Glugur Kota; rating 4,7 dari 4.584 penilaian; status saat screenshot tutup dan buka besok; jarak 6,6 km. \\ \hline
        TO-GF-02 & Ringkasan ulasan agregat & Ringkasan ulasan menonjolkan rasa enak, porsi besar, bahan segar, kemasan rapi, menu beragam, dan harga terjangkau. \\ \hline
        TO-GF-03 & Promo dan voucher & Terdapat 19 promo/voucher aktif yang terlihat, meliputi diskon persentase, potongan nominal, diskon ongkir, bank, dan \textit{group order}. \\ \hline
        TO-GF-04 & Menu terlaris & Menu yang terlihat berlabel terlaris adalah Nasi Padang Rendang Sapi, Nasi Padang Ayam Goreng Bumbu, dan Nasi Padang Ayam Rendang. \\ \hline
        TO-GF-05 & Rentang harga & Rentang harga observasi Rp1.100--Rp275.000; item tersedia Rp1.100--Rp72.000; paket nasi Padang tersedia Rp27.500--Rp72.000. \\ \hline
        TO-GF-06 & Ketersediaan menu & Terdapat 6 menu atau produk yang terlihat berstatus \textit{Nggak tersedia}, termasuk menu premium kepala ikan dan produk sambal kemasan. \\ \hline
        TO-GF-07 & Ulasan positif & Tema positif utama meliputi rasa enak, porsi besar, kesegaran bahan, kebersihan, pemenuhan request, pengemasan, dan pelanggan rutin. \\ \hline
        TO-GF-08 & Ulasan negatif/netral & Tema negatif/netral meliputi salah item, item kurang lengkap, kuah/add-on tidak diberikan, persepsi harga/porsi, kualitas/kesegaran, koordinasi staf, dan risiko pengalaman pascakonsumsi. \\ \hline
    \end{longtable}
\endgroup

Observasi kompetitor pada 19 dan 20 Juni 2026 memperlihatkan lima pembanding utama. Perbandingan reputasi digital dan harga terlihat disajikan pada Tabel~\ref{tab:perbandingan-reputasi-digital-kompetitor-bab4} dan Tabel~\ref{tab:perbandingan-rentang-harga-kompetitor-bab4}.

\begingroup
    \scriptsize
    \setlength{\tabcolsep}{2pt}
    \renewcommand{\arraystretch}{1.02}
    \begin{longtable}{|Z{2.45cm}|Z{1.05cm}|Z{1.55cm}|Z{1.15cm}|L{6.45cm}|}
        \caption{Perbandingan Reputasi Digital Nasi Gerilya dan Kompetitor}
        \label{tab:perbandingan-reputasi-digital-kompetitor-bab4}\\
        \toprule
        \textbf{Merchant} & \textbf{Rating} & \textbf{Penilaian} & \textbf{Jarak} & \textbf{Makna} \\ \hline
        \midrule
        \endfirsthead
        \multicolumn{5}{@{}l@{}}{\footnotesize\textit{Lanjutan Tabel \thetable}}\\
        \toprule
        \textbf{Merchant} & \textbf{Rating} & \textbf{Penilaian} & \textbf{Jarak} & \textbf{Makna} \\ \hline
        \midrule
        \endhead
        \bottomrule
        \multicolumn{5}{r}{\footnotesize\textit{Bersambung ke halaman berikutnya}}\\
        \endfoot
        \bottomrule
        \multicolumn{5}{@{}L{\textwidth}@{}}{\footnotesize Sumber: Diolah peneliti dari observasi GrabFood Nasi Gerilya dan kompetitor (2026).}\\
        \endlastfoot
        Nasi Gerilya - Glugur Kota & 4,7 & 4.584 & 6,6 km & Rating kuat, tetapi jumlah penilaian masih di bawah Pondok Gurih, Garuda, dan Istana Krakatau. \\ \hline
        Restoran Paripurna & 4,7 & 1.264 & 3,0 km & Rating setara dan lebih dekat pada titik observasi. \\ \hline
        Rumah Makan Dendeng Batokok & 4,8 & 1.479 & 7,3 km & Rating lebih tinggi dan menjadi pembanding harga bawah. \\ \hline
        Rumah Makan Istana Krakatau & 4,6 & 6.256 & 8,4 km & Rating lebih rendah, tetapi jumlah penilaian dan menu premium lebih besar. \\ \hline
        RM. Pondok Gurih & 4,9 & 56.157 & 5,0 km & Rating dan jumlah penilaian paling tinggi serta memiliki penanda reputasi kuat. \\ \hline
        Restoran Garuda & 4,6 & 17.068 & 6,7 km & Jumlah penilaian jauh lebih besar dan jaraknya hampir sama dengan Nasi Gerilya. \\ \hline
    \end{longtable}
\endgroup

\begingroup
    \scriptsize
    \setlength{\tabcolsep}{2pt}
    \renewcommand{\arraystretch}{1.02}
    \begin{longtable}{|Z{2.45cm}|Z{2.75cm}|Z{3.1cm}|L{4.45cm}|}
        \caption{Perbandingan Rentang Harga Terlihat pada GrabFood}
        \label{tab:perbandingan-rentang-harga-kompetitor-bab4}\\
        \toprule
        \textbf{Merchant} & \textbf{Rentang Harga Terlihat} & \textbf{Rentang Menu Utama/Paket} & \textbf{Makna} \\ \hline
        \midrule
        \endfirsthead
        \multicolumn{4}{@{}l@{}}{\footnotesize\textit{Lanjutan Tabel \thetable}}\\
        \toprule
        \textbf{Merchant} & \textbf{Rentang Harga Terlihat} & \textbf{Rentang Menu Utama/Paket} & \textbf{Makna} \\ \hline
        \midrule
        \endhead
        \bottomrule
        \multicolumn{4}{r}{\footnotesize\textit{Bersambung ke halaman berikutnya}}\\
        \endfoot
        \bottomrule
        \multicolumn{4}{@{}L{\textwidth}@{}}{\footnotesize Sumber: Diolah peneliti dari observasi GrabFood Nasi Gerilya dan kompetitor (2026).}\\
        \endlastfoot
        Nasi Gerilya & Keseluruhan Rp1.100--Rp275.000; item tersedia Rp1.100--Rp72.000 & Paket nasi Padang tersedia Rp27.500--Rp72.000 & Rentang luas, tetapi pelanggan tetap sensitif terhadap harga akhir, promo, dan ongkir. \\ \hline
        Paripurna & Rp5.000--Rp66.000 & Nasi bungkus utama Rp24.300--Rp66.000 & Dekat dengan rentang paket Nasi Gerilya. \\ \hline
        Dendeng Batokok & Rp5.000--Rp36.000 & Nasi/lauk Rp7.000--Rp36.000 & Batas harga atas lebih rendah, sehingga menjadi ancaman untuk pelanggan sensitif harga. \\ \hline
        Istana Krakatau & Rp9.000--Rp250.000 & Nasi bungkus/kotak Rp22.500--Rp83.000 & Memiliki spektrum premium dan menu ikan yang kuat. \\ \hline
        Pondok Gurih & Rp6.750--Rp246.750 & Nasi/lauk Rp23.000--Rp75.000 & Reputasi tinggi dengan harga luas. \\ \hline
        Garuda & Rp14.300--Rp265.851 & Nasi bungkus/kotak Rp19.048--Rp92.063 & Akumulasi penilaian besar dan rentang premium. \\ \hline
    \end{longtable}
\endgroup

\section{Analisis Strategi Pemasaran Rumah Makan Nasi Gerilya}

\subsection{Analisis STP}
Analisis STP digunakan untuk membaca pasar secara lebih terarah karena strategi pemasaran perlu dimulai dari pemilihan segmen, penentuan target, dan pembentukan posisi yang jelas dalam benak pelanggan \parencite{kotler2023principlesmarketing,kotler2021marketingmanagement}. Pada konteks platform pesan-antar, posisi tersebut tidak hanya dibentuk oleh rasa makanan, tetapi juga oleh foto, rating, ulasan, harga akhir, promo, dan kemudahan informasi yang dilihat pelanggan sebelum melakukan checkout \parencite{chaffey2022digitalmarketing,sophia2025rendanggadih}. Berdasarkan temuan lapangan, Nasi Gerilya melayani pelanggan GrabFood yang mencari makanan berat, praktis, mengenyangkan, dan dapat dipercaya melalui tampilan digital. Ringkasan STP disajikan pada Tabel~\ref{tab:analisis-stp-bab4}.

\begingroup
    \footnotesize
    \setlength{\tabcolsep}{2pt}
    \renewcommand{\arraystretch}{1.05}
    \begin{longtable}{|Z{2.25cm}|L{4.05cm}|L{2.45cm}|L{4.25cm}|}
        \caption{Analisis STP Nasi Gerilya pada GrabFood}
        \label{tab:analisis-stp-bab4}\\
        \toprule
        \textbf{Komponen} & \textbf{Temuan Lapangan} & \textbf{Bukti Data} & \textbf{Implikasi Strategi} \\ \hline
        \midrule
        \endfirsthead
        \multicolumn{4}{@{}l@{}}{\footnotesize\textit{Lanjutan Tabel \thetable}}\\
        \toprule
        \textbf{Komponen} & \textbf{Temuan Lapangan} & \textbf{Bukti Data} & \textbf{Implikasi Strategi} \\ \hline
        \midrule
        \endhead
        \bottomrule
        \multicolumn{4}{r}{\footnotesize\textit{Bersambung ke halaman berikutnya}}\\
        \endfoot
        \bottomrule
        \multicolumn{4}{@{}L{\textwidth}@{}}{\footnotesize Sumber: Diolah peneliti berdasarkan hasil penelitian utama (2026).}\\
        \endlastfoot
        Segmentasi & Pelanggan GrabFood terdiri atas pekerja, mahasiswa, pelanggan sekitar, dan pelanggan jam makan siang yang membutuhkan makanan berat saat lelah, sibuk, panas, malam, atau tidak sempat keluar. & TS-PL-01, TS-PL-02, W-BC-07 & Segmen utama perlu diarahkan pada makan siang pekerja, makan malam mahasiswa, pelanggan sekitar yang sensitif terhadap ongkir, serta pelanggan lintas kanal yang membutuhkan layanan cepat. \\ \hline
        \textit{Targeting} & Target paling potensial adalah pelanggan yang mencari nasi Padang praktis, porsi besar, menu lauk kuat, rating tinggi, dan promo yang menurunkan harga akhir tanpa mengabaikan margin. & TS-PL-03, TS-PL-04, TS-PL-06, TO-GF-03, W-BC-07, W-BC-10 & Buat paket menu untuk pekerja/kantor, anak kos, pelanggan promo, dan pelanggan makan siang dengan tetap menghitung HPP dan kesiapan stok. \\ \hline
        \textit{Positioning} & Nasi Gerilya paling kuat diposisikan sebagai nasi Padang Medan yang praktis, cepat, mengenyangkan, bercita rasa lauk kuat, dan didukung tampilan digital yang meyakinkan. & TO-GF-01, TO-GF-02, TW-KR-05, TS-PL-03, W-BC-08, W-BC-09 & Positioning perlu ditegaskan melalui foto menu, deskripsi isi paket, highlight dendeng/ayam pop/ayam goreng, menu digital, serta janji ketelitian pesanan. \\ \hline
    \end{longtable}
\endgroup

\subsection{Analisis Bauran Pemasaran 7P}
Analisis bauran pemasaran 7P digunakan karena usaha kuliner pada platform digital tidak hanya menjual produk, tetapi juga mengelola harga, kanal distribusi, promosi, orang, proses, dan bukti fisik/digital yang membentuk pengalaman pelanggan \parencite{kotler2023principlesmarketing,wirtz2024servicesmarketing}. Dalam pemasaran digital, bukti seperti foto, rating, deskripsi menu, ulasan, dan tampilan toko menjadi bagian dari proses evaluasi pelanggan sebelum membeli \parencite{chaffey2022digitalmarketing}. Hasil analisis 7P menunjukkan bahwa produk dan bukti digital Nasi Gerilya sudah cukup kuat, tetapi proses layanan menjadi titik paling penting untuk diperbaiki. Ringkasan 7P disajikan pada Tabel~\ref{tab:analisis-7p-bab4}.

\begingroup
    \scriptsize
    \setlength{\tabcolsep}{2pt}
    \renewcommand{\arraystretch}{1.02}
    \begin{longtable}{|Z{1.8cm}|L{5.05cm}|L{2.95cm}|L{3.05cm}|}
        \caption{Analisis Bauran Pemasaran 7P Nasi Gerilya}
        \label{tab:analisis-7p-bab4}\\
        \toprule
        \textbf{Unsur 7P} & \textbf{Temuan} & \textbf{Kode Data} & \textbf{Implikasi} \\ \hline
        \midrule
        \endfirsthead
        \multicolumn{4}{@{}l@{}}{\footnotesize\textit{Lanjutan Tabel \thetable}}\\
        \toprule
        \textbf{Unsur 7P} & \textbf{Temuan} & \textbf{Kode Data} & \textbf{Implikasi} \\ \hline
        \midrule
        \endhead
        \bottomrule
        \multicolumn{4}{r}{\footnotesize\textit{Bersambung ke halaman berikutnya}}\\
        \endfoot
        \bottomrule
        \multicolumn{4}{@{}L{\textwidth}@{}}{\footnotesize Sumber: Diolah peneliti berdasarkan hasil penelitian utama (2026).}\\
        \endlastfoot
        \textit{Product} & Menu unggulan kuat pada dendeng, ayam pop, ayam goreng, serta porsi besar. Ayam pop juga dinilai layak menjadi hero menu online, sementara konsistensi rasa tetap perlu dijaga melalui standar resep dan double-check. & TW-KR-05, TW-KR-06, TS-PL-03, TS-PL-04, W-BC-09 & Pertahankan menu hero dan tambah paket berbasis segmen, khususnya paket ayam pop untuk makan siang/kantor. \\ \hline
        \textit{Price} & Harga dinilai cukup mahal oleh pelanggan, terutama ketika tidak ada promo atau ongkir murah. Promo membantu menarik pembeli, tetapi perlu dihitung dengan HPP dan margin. & TW-AD-09, TS-PL-05, TS-PL-06, W-BC-10, W-BC-17 & Harga perlu dikemas dalam paket hemat, bundling, komunikasi nilai porsi, dan batas promo yang tidak mengganggu margin. \\ \hline
        \textit{Place} & GrabFood menjadi kanal utama yang memudahkan pemesanan, tetapi Nasi Gerilya juga melayani dine in dan take away sehingga antrean lintas kanal memengaruhi kecepatan layanan. & TO-GF-01, TS-PL-01, TS-PL-10, W-BC-05, W-BC-06 & Update stok, estimasi waktu, dan sistem antrean perlu dijaga agar pengalaman kanal digital dan offline tetap terkendali. \\ \hline
        \textit{Promotion} & Promo mendorong pesanan, trial, dan pembelian ulang, tetapi dapat membuat order membeludak jika dapur tidak siap. Stamp digital membuka peluang loyalitas langsung di luar promo platform. & TO-GF-03, TW-AD-10, TS-PL-06, W-BC-05, W-BC-12, W-BC-16 & Promo harus disinkronkan dengan stok/dapur, sedangkan stamp digital dapat digunakan untuk mengurangi ketergantungan pada voucher platform. \\ \hline
        \textit{People} & Kasir, checker, display, dan dapur sudah memiliki peran, tetapi kondisi ramai menuntut koordinasi, keramahan, dan kemampuan memakai sistem digital secara konsisten. & TW-KR-01, TW-KR-09, TW-KR-10, TW-AD-01, W-BC-01, W-BC-03 & Perlu briefing SOP, pembagian tugas saat peak hour, dan pembiasaan penggunaan sistem antrean/menu digital. \\ \hline
        \textit{Process} & Alur pesanan online sudah ada, tetapi rawan item tertinggal, nomor mirip, catatan terlewat, driver menunggu, menu habis, dan antrean tiga jalur yang menumpuk. & TW-AD-04, TW-AD-05, TW-AD-06, TS-PL-09, TS-PL-10, W-BC-05, W-BC-19, W-BC-20 & Checklist packing, update stok real-time, dan sistem antrean menjadi prioritas proses. \\ \hline
        \textit{Physical evidence} & Rating, foto, ulasan, kemasan, dan desain visual cukup kuat, tetapi pelanggan meminta detail isi paket, kuah/sambal, segel, label, informasi stok, dan pengalaman digital yang lebih mandiri. & TO-GF-01, TO-GF-02, TS-PL-07, TS-PL-08, W-BC-11 & Perbaiki foto, deskripsi, label, tampilan paket unggulan di GrabFood, serta konsistensi menu digital. \\ \hline
    \end{longtable}
\endgroup

\section{Analisis Pelanggan, Ulasan, dan Indikator Penjualan}
Pengalaman pelanggan menunjukkan bahwa keputusan pembelian tidak ditentukan oleh satu faktor. Dalam teori perilaku pelanggan dan pemasaran jasa, nilai yang dirasakan pelanggan terbentuk dari manfaat produk, kemudahan layanan, risiko yang dirasakan, serta pengalaman selama proses pembelian \parencite{kotler2023principlesmarketing,wirtz2024servicesmarketing}. Pada platform digital, evaluasi tersebut terjadi secara cepat karena pelanggan dapat membandingkan rating, foto, promo, ongkir, dan ulasan antarrestoran sebelum memutuskan pembelian \parencite{chaffey2022digitalmarketing}. Pelanggan Nasi Gerilya membeli karena kombinasi rasa, porsi, menu unggulan, rating, foto, jarak, promo, dan ongkir. Namun, pelanggan juga mudah menunda atau berpindah ketika harga akhir terlalu tinggi, voucher tidak tersedia, stok tidak jelas, atau pesanan kurang lengkap. Ringkasan hubungan pengalaman pelanggan dengan penjualan disajikan pada Tabel~\ref{tab:analisis-pelanggan-ulasan}.

Data kinerja pra-survei memperlihatkan bahwa isu penjualan tidak cukup dibaca hanya dari jumlah transaksi. Pola awal menunjukkan penjualan bulanan sempat mencapai puncak pada Maret 2025 lalu cenderung melemah sampai November 2025, sedangkan indeks Nasi Gerilya terhadap kompetitor turun dari 100,0 menjadi 70,4 pada Juli--Oktober 2025. Pada saat yang sama, AOV dan kunjungan naik pada sebagian periode, sehingga masalah yang perlu dibaca bukan hanya traffic, tetapi juga konversi, repeat order, persepsi harga, promo, stok, dan pengalaman layanan. Oleh karena itu, data kinerja digunakan sebagai konteks urgensi dan indikator evaluasi strategi, bukan sebagai pembuktian kuantitatif dampak strategi.

\begingroup
    \footnotesize
    \setlength{\tabcolsep}{2pt}
    \renewcommand{\arraystretch}{1.05}
    \begin{longtable}{|Z{2.35cm}|L{3.85cm}|L{3.35cm}|L{3.35cm}|}
        \caption{Analisis Pelanggan dan Hubungannya dengan Penjualan}
        \label{tab:analisis-pelanggan-ulasan}\\
        \toprule
        \textbf{Tema} & \textbf{Data Pendukung} & \textbf{Makna Temuan} & \textbf{Hubungan dengan Penjualan} \\ \hline
        \midrule
        \endfirsthead
        \multicolumn{4}{@{}l@{}}{\footnotesize\textit{Lanjutan Tabel \thetable}}\\
        \toprule
        \textbf{Tema} & \textbf{Data Pendukung} & \textbf{Makna Temuan} & \textbf{Hubungan dengan Penjualan} \\ \hline
        \midrule
        \endhead
        \bottomrule
        \multicolumn{4}{r}{\footnotesize\textit{Bersambung ke halaman berikutnya}}\\
        \endfoot
        \bottomrule
        \multicolumn{4}{@{}L{\textwidth}@{}}{\footnotesize Sumber: Diolah peneliti berdasarkan angket pelanggan, ulasan, dan observasi kompetitor (2026).}\\
        \endlastfoot
        Alasan membeli & Dendeng, ayam pop, ayam goreng, porsi besar, rating, foto, promo, dan jarak dekat. & Menu hero dan bukti digital membentuk minat awal. & Meningkatkan peluang trial dan order pertama. \\ \hline
        Alasan membeli ulang & Rasa konsisten, porsi, promo, rating, dan respons keluhan. & Repeat order bergantung pada produk dan proses. & Menjaga pelanggan lama dan mengurangi penurunan repeat order. \\ \hline
        Keluhan utama & Item kurang, catatan khusus, stok, makanan dingin, dan harga akhir. & Keluhan menyentuh proses, harga, dan bukti layanan. & Jika tidak diperbaiki, keluhan dapat menurunkan rating dan pembelian ulang. \\ \hline
        Persepsi harga/promo & Harga dianggap mahal, tetapi promo dan ongkir murah membuatnya diterima. & Promo menjadi penyeimbang persepsi harga. & Mendorong order, tetapi perlu dijaga agar tidak hanya bergantung pada voucher. \\ \hline
        Peluang upselling & Pra-survei menunjukkan add-on, lauk tambahan, sambal kemasan, kuah/minuman, dan menu premium tersedia sebagai variasi penawaran. & Product lineup tidak hanya berfungsi sebagai variasi menu, tetapi juga peluang peningkatan AOV. & Mendukung bundling menu hero, sambal/kuah/minuman, dan paket segmen. \\ \hline
        Indikasi kinerja pra-survei & Penjualan bulanan dan indeks relatif melemah, sementara AOV/kunjungan naik pada sebagian periode. & Ada indikasi bahwa persoalan tidak hanya pada minat awal, tetapi juga konversi, repeat order, dan pengalaman layanan. & Strategi diarahkan pada perbaikan proses dan loyalitas, bukan sekadar menambah promo. \\ \hline
        Kompetitor & Dendeng Batokok, Paripurna, Pondok Gurih, Garuda, dan Istana Krakatau menjadi pembanding. & Pelanggan aktif membandingkan rasa, harga, porsi, dan promo. & Nasi Gerilya perlu memperjelas pembeda agar tidak mudah diganti. \\ \hline
    \end{longtable}
\endgroup

\section{Identifikasi Faktor Internal dan Eksternal}
Faktor internal dan eksternal dirumuskan dari koding temuan lapangan, lalu diperkuat oleh kode pra-survei jika temanya berulang pada sumber utama. Dalam kerangka manajemen strategis, kekuatan dan kelemahan dibaca sebagai kondisi internal yang relatif dapat dikendalikan organisasi, sedangkan peluang dan ancaman berasal dari lingkungan eksternal yang perlu direspons melalui strategi \parencite{david2023strategicmanagement,gurel2017swot}. Karena penelitian ini menggunakan SWOT sebagai alat sintesis, faktor yang dimasukkan adalah faktor yang memiliki hubungan langsung dengan daya saing, keputusan pembelian, repeat order, atau risiko penurunan penjualan \parencite{rangkuti2016analisisSWOT}. Faktor internal terdiri dari kekuatan dan kelemahan yang relatif dapat dikendalikan usaha. Faktor eksternal terdiri dari peluang dan ancaman yang muncul dari pelanggan, platform, kompetitor, dan perubahan perilaku pembelian.

\begingroup
    \scriptsize
    \setlength{\tabcolsep}{2pt}
    \renewcommand{\arraystretch}{1.02}
    \begin{longtable}{|Z{0.9cm}|L{5.3cm}|L{3.65cm}|L{2.9cm}|}
        \caption{Identifikasi Faktor Internal dan Eksternal}
        \label{tab:faktor-internal-eksternal-bab4}\\
        \toprule
        \textbf{Kode} & \textbf{Faktor} & \textbf{Bukti Data} & \textbf{Dampak terhadap Penjualan} \\ \hline
        \midrule
        \endfirsthead
        \multicolumn{4}{@{}l@{}}{\footnotesize\textit{Lanjutan Tabel \thetable}}\\
        \toprule
        \textbf{Kode} & \textbf{Faktor} & \textbf{Bukti Data} & \textbf{Dampak terhadap Penjualan} \\ \hline
        \midrule
        \endhead
        \bottomrule
        \multicolumn{4}{r}{\footnotesize\textit{Bersambung ke halaman berikutnya}}\\
        \endfoot
        \bottomrule
        \multicolumn{4}{@{}L{\textwidth}@{}}{\footnotesize Sumber: Diolah peneliti berdasarkan koding temuan lapangan (2026).}\\
        \endlastfoot
        S1 & Rasa dan menu unggulan kuat, terutama dendeng, ayam pop, ayam goreng, dan variasi lauk. & TO-GF-02, TO-GF-04, TW-KR-05, TS-PL-03, W-BC-08, W-BC-09, PS-WPM-011, PS-PROD-003, PS-PROD-014 & Menarik trial dan repeat order. \\ \hline
        S2 & Porsi besar dan kemasan relatif rapi memberi nilai pelanggan. & TO-GF-07, TS-PL-04, TS-PL-08 & Meningkatkan persepsi value. \\ \hline
        S3 & Alur operasional online sudah memiliki checker, print bill, nomor antrean, update lauk, koordinasi display/dapur, dan arah digitalisasi internal. & TW-AD-03, TW-AD-04, TW-KR-01, W-BC-01, W-BC-03, W-BC-19 & Menjadi dasar layanan online yang dapat diperkuat. \\ \hline
        S4 & Reputasi digital cukup kuat dengan rating 4,7, 4.584 penilaian, promo aktif, serta peluang penguatan tampilan/menu digital. & TO-GF-01, TO-GF-03, W-BC-11, PS-GF-001 & Membantu keyakinan pelanggan baru. \\ \hline
        S5 & Respons keluhan sudah ada melalui pengiriman ulang item kurang dan permintaan maaf/komplimen. & TW-AD-11, TS-PL-09 & Mengurangi dampak negatif kesalahan pesanan. \\ \hline
        W1 & Akurasi packing rawan saat ramai, terutama item terpisah/berkuah, catatan khusus, dan nomor GrabFood mirip. & TO-GF-08, TW-AD-05, TW-AD-06, TS-PL-09, PS-WPM-020 & Menurunkan kepuasan dan dapat memicu ulasan negatif. \\ \hline
        W2 & Kecepatan layanan terhambat saat peak hour karena antrean dine in, ojek online, take away, pesanan mendadak, SDM, packing, dan driver menunggu. & TW-KR-03, TW-KR-04, TW-AD-07, W-BC-05, W-BC-15, PS-WPM-018 & Memengaruhi estimasi waktu dan repeat order. \\ \hline
        W3 & Update stok dan kesiapan menu belum sepenuhnya stabil. & TO-GF-06, TW-KR-08, TS-PL-10, PS-WPM-017 & Membuat pelanggan batal atau berpindah menu/restoran. \\ \hline
        W4 & Persepsi harga di GrabFood cenderung mahal dan bergantung pada promo/ongkir, sementara promo tetap harus mempertimbangkan HPP dan margin. & TW-AD-09, TS-PL-05, TS-PL-06, W-BC-10, PS-WPM-014, PS-WPM-015 & Menurunkan minat beli tanpa voucher dan dapat menekan margin bila promo tidak dihitung. \\ \hline
        W5 & Detail menu, isi paket, kuah/sambal, segel, label, batas request, dan data pelanggan langsung masih perlu dipertegas. & TW-KR-07, TW-AD-12, TS-PL-07, TS-PL-08, W-BC-16 & Membuat pelanggan ragu atau kecewa jika ekspektasi tidak sesuai, serta membuat loyalitas terlalu bergantung pada platform. \\ \hline
        O1 & Pekerja, mahasiswa, pelanggan sekitar, dan pelanggan jam makan siang membutuhkan makanan berat praktis melalui GrabFood. & TS-PL-01, TS-PL-02, W-BC-07, W-BC-17 & Membuka peluang paket makan siang/malam. \\ \hline
        O2 & Promo, voucher, ongkir murah, paket segmen, upselling, dan stamp digital dapat mendorong trial, repeat order, serta loyalitas langsung. & TO-GF-03, TW-AD-10, TS-PL-06, W-BC-05, W-BC-12, W-BC-16, PS-WPM-022, PS-GF-031--PS-GF-037 & Meningkatkan volume pesanan saat dikendalikan dengan stok dan memperkuat retensi pelanggan. \\ \hline
        O3 & Tampilan digital, foto, rating, ulasan, deskripsi, dan menu digital dapat diperkuat untuk menaikkan rasa percaya. & TO-GF-01, TO-GF-02, TS-PL-07, W-BC-11 & Membantu konversi dari lihat menu menjadi checkout dan memudahkan pemesanan mandiri. \\ \hline
        O4 & Kompetitor juga menghadapi masalah stok, item tertinggal, porsi/rasa tidak konsisten, dan driver menunggu. & TW-KP-PK-06--TW-KP-PK-08, TW-KP-IK-06--TW-KP-IK-07 & Nasi Gerilya dapat membedakan diri lewat ketelitian proses. \\ \hline
        O5 & Media sosial, rekomendasi teman, ulasan pelanggan, dan database pelanggan dapat memperkuat awareness serta hubungan pelanggan di luar GrabFood. & TW-AD-10, TS-PL-07, W-BC-12, W-BC-16 & Membuka sumber pelanggan baru di luar pencarian GrabFood dan mengurangi ketergantungan pada trafik platform. \\ \hline
        T1 & Kompetitor memiliki reputasi digital kuat, terutama Pondok Gurih, Garuda, Istana Krakatau, dan kompetitor dekat yang muncul pada data pra-survei. & TO-KP-01, TO-KP-02, TO-KP-07, PS-WPM-024, PS-KIN-035, PS-KIN-036 & Menekan posisi Nasi Gerilya pada tahap perbandingan pelanggan. \\ \hline
        T2 & Pelanggan mudah membandingkan harga, promo, ongkir, dan menu dengan kompetitor. & TO-KP-03, TO-KP-06, TS-PL-12 & Pelanggan dapat berpindah bila harga akhir lebih menarik di tempat lain. \\ \hline
        T3 & Biaya tambahan, ongkir, biaya platform, kenaikan bahan baku, dan ketergantungan voucher menurunkan minat beli tanpa promo. & TS-PL-05, TS-PL-06, TS-PL-10, W-BC-17, PS-WPM-014, PS-WPM-024 & Membuat permintaan tidak stabil dan menekan ruang margin. \\ \hline
        T4 & Human error, makanan dingin, stok habis, antrean menumpuk, dan item kurang dapat memunculkan ulasan negatif. & TO-GF-08, TS-PL-10, TS-PL-11, W-BC-15 & Mengganggu rating dan loyalitas pelanggan. \\ \hline
        T5 & Kompetitor menonjolkan rasa, pelayanan cepat, sambal, porsi, dan menu spesifik. & TW-KP-PK-04, TW-KP-IK-04, TW-KP-IK-09 & Memperkuat tekanan diferensiasi produk dan layanan. \\ \hline
    \end{longtable}
\endgroup

\section{Triangulasi Data}
Triangulasi dilakukan dengan membandingkan temuan dari pegawai, intern \textit{business consultant}, pelanggan, kompetitor, observasi platform, dan dokumentasi kinerja. Dalam penelitian kualitatif, triangulasi dan proses kondensasi data membantu peneliti memeriksa konsistensi makna antarsumber sebelum menyusun kesimpulan tematik \parencite{creswell2018qualitativeinquiry,miles2020qualitativedata}. Oleh karena itu, faktor yang digunakan dalam IFAS, EFAS, dan SWOT tidak diambil dari satu kutipan tunggal, tetapi dari pola yang muncul berulang atau memiliki dampak strategis langsung. Hasil triangulasi disajikan pada Tabel~\ref{tab:triangulasi-bab4}.

\begingroup
    \scriptsize
    \setlength{\tabcolsep}{2pt}
    \renewcommand{\arraystretch}{1.02}
    \begin{longtable}{|Z{2.0cm}|L{3.35cm}|L{2.45cm}|L{2.05cm}|Z{2.25cm}|}
        \caption{Triangulasi Temuan Utama}
        \label{tab:triangulasi-bab4}\\
        \toprule
        \textbf{Temuan} & \textbf{Sumber} & \textbf{Teknik} & \textbf{Waktu} & \textbf{Keputusan} \\ \hline
        \midrule
        \endfirsthead
        \multicolumn{5}{@{}l@{}}{\footnotesize\textit{Lanjutan Tabel \thetable}}\\
        \toprule
        \textbf{Temuan} & \textbf{Sumber} & \textbf{Teknik} & \textbf{Waktu} & \textbf{Keputusan} \\ \hline
        \midrule
        \endhead
        \bottomrule
        \multicolumn{5}{r}{\footnotesize\textit{Bersambung ke halaman berikutnya}}\\
        \endfoot
        \bottomrule
        \multicolumn{5}{@{}L{\textwidth}@{}}{\footnotesize Sumber: Diolah peneliti berdasarkan triangulasi data penelitian utama dan pra-survei (2026).}\\
        \endlastfoot
        Rasa dan menu unggulan menjadi kekuatan utama. & Pegawai, pelanggan, ulasan, intern \textit{business consultant} & Wawancara, angket, observasi & Juni--Juli 2026 & Valid \\ \hline
        Promo dan ongkir memengaruhi keputusan beli. & Pelanggan, pegawai, observasi promo & Angket, wawancara, observasi & Juni--Juli 2026 & Valid \\ \hline
        Akurasi packing dan antrean lintas kanal menjadi titik lemah proses. & Pegawai, pelanggan, ulasan, intern \textit{business consultant} & Wawancara, angket, observasi & Juni--Juli 2026 & Valid \\ \hline
        Update stok dan menu habis perlu diperbaiki. & Pegawai, pelanggan, observasi & Wawancara, angket, observasi & Juni--Juli 2026 & Valid dengan catatan \\ \hline
        Masalah stok, antrean, promo, sensitivitas harga, dan repeat order sudah tampak sejak pra-survei dan muncul kembali pada data lapangan. & Wawancara awal pemilik, dokumentasi kinerja, pegawai, pelanggan, intern \textit{business consultant} & Wawancara awal, dokumentasi, wawancara, angket & April--Juli 2026 & Valid sebagai pola berulang \\ \hline
        Kompetitor kuat dari reputasi, rasa, dan layanan cepat. & Observasi kompetitor, wawancara kompetitor, pelanggan & Observasi, wawancara, angket & Juni--Juli 2026 & Valid \\ \hline
        Ketergantungan pelanggan pada voucher dapat menekan pembelian tanpa promo. & Pelanggan, kasir, intern \textit{business consultant} & Angket dan wawancara & Juli 2026 & Valid \\ \hline
        Sistem antrean, menu digital, dan stamp digital relevan untuk memperkuat pengalaman layanan dan loyalitas. & Intern \textit{business consultant}, pegawai, pelanggan & Wawancara dan angket & Juli 2026 & Valid sebagai rekomendasi bertahap \\ \hline
        Promo perlu disesuaikan dengan HPP, margin, stok, dan kapasitas dapur. & Wawancara awal pemilik, pegawai, intern \textit{business consultant}, pelanggan & Wawancara dan angket & April--Juli 2026 & Valid sebagai pertimbangan implementasi \\ \hline
    \end{longtable}
\endgroup

\section{Matriks IFAS dan EFAS}
Penyusunan IFAS dan EFAS mengikuti prosedur Bab III. Dalam matriks IFAS dan EFAS, bobot menunjukkan tingkat kepentingan relatif setiap faktor terhadap keberhasilan strategi, sedangkan rating menunjukkan kondisi respons usaha terhadap faktor tersebut \parencite{rangkuti2016analisisSWOT,david2023strategicmanagement}. Bobot ditentukan berdasarkan tiga pertimbangan: kekuatan bukti, frekuensi kemunculan tema antarsumber, dan dampak faktor terhadap penjualan atau pembelian ulang. Faktor yang muncul pada banyak sumber dan langsung memengaruhi trial, repeat order, rating, komplain, atau margin diberi bobot lebih tinggi. Rating internal menunjukkan tingkat kekuatan atau kelemahan, sedangkan rating eksternal menunjukkan kemampuan respons usaha terhadap peluang atau ancaman. Dengan demikian, bobot dan rating pada matriks ini dibaca sebagai penilaian analitis berbasis triangulasi sumber untuk menentukan prioritas strategi.

Justifikasi bobot disajikan secara rinci pada Tabel~\ref{tab:justifikasi-bobot-ifas-efas-bab4}. Penjelasan ini menunjukkan hubungan antara tingkat kepentingan faktor, pola temuan lapangan, dan dampak strategis masing-masing faktor terhadap penjualan \parencite{hendrizon2024madayacoffee,rangkuti2016analisisSWOT}.

\begingroup
    \scriptsize
    \setlength{\tabcolsep}{2pt}
    \renewcommand{\arraystretch}{1.02}
    \begin{longtable}{|Z{1.05cm}|Z{0.8cm}|C{0.9cm}|L{9.65cm}|}
        \caption{Justifikasi Bobot IFAS dan EFAS}
        \label{tab:justifikasi-bobot-ifas-efas-bab4}\\
        \toprule
        \textbf{Matriks} & \textbf{Kode} & \textbf{Bobot} & \textbf{Justifikasi} \\ \hline
        \midrule
        \endfirsthead
        \multicolumn{4}{@{}l@{}}{\footnotesize\textit{Lanjutan Tabel \thetable}}\\
        \toprule
        \textbf{Matriks} & \textbf{Kode} & \textbf{Bobot} & \textbf{Justifikasi} \\ \hline
        \midrule
        \endhead
        \bottomrule
        \multicolumn{4}{r}{\footnotesize\textit{Bersambung ke halaman berikutnya}}\\
        \endfoot
        \bottomrule
        \multicolumn{4}{@{}L{\textwidth}@{}}{\footnotesize Sumber: Diolah peneliti berdasarkan triangulasi temuan lapangan, pemetaan pra-survei, dan konsep IFAS/EFAS \parencite{rangkuti2016analisisSWOT,david2023strategicmanagement,hendrizon2024madayacoffee}.}\\
        \endlastfoot
        IFAS & S1 & 0,14 & Bobot tertinggi pada kekuatan diberikan karena rasa dan menu unggulan muncul pada pegawai, pelanggan, intern \textit{business consultant}, observasi, dan pra-survei. Faktor ini berhubungan langsung dengan trial, repeat order, dan positioning Nasi Gerilya sebagai nasi Padang Medan yang mengenyangkan dan bercita rasa kuat. \\ \hline
        IFAS & S2 & 0,10 & Porsi besar dan kemasan rapi penting karena membentuk persepsi nilai pelanggan. Bobotnya tinggi, tetapi tidak setinggi S1 karena porsi dan kemasan mendukung keputusan beli setelah daya tarik menu utama terbentuk. \\ \hline
        IFAS & S3 & 0,11 & Alur operasional online, checker, bill, nomor antrean, update lauk, dan arah digitalisasi internal menjadi dasar untuk mengeksekusi layanan. Bobotnya lebih tinggi daripada S4 dan S5 karena proses menentukan apakah pesanan digital dapat dilayani konsisten, tetapi belum diberi bobot tertinggi karena sistemnya masih perlu dimatangkan. \\ \hline
        IFAS & S4 & 0,09 & Reputasi digital dan promo aktif membantu pelanggan baru percaya sebelum membeli. Bobotnya sedang karena rating, foto, dan promo penting pada tahap konversi, tetapi efeknya akan melemah jika proses packing, stok, dan antrean tidak stabil. \\ \hline
        IFAS & S5 & 0,07 & Respons keluhan menjadi kekuatan pendukung karena dapat mengurangi dampak negatif kesalahan layanan. Bobotnya paling kecil di antara kekuatan karena sifatnya reaktif dan muncul pada lebih sedikit data dibandingkan rasa, porsi, dan proses. \\ \hline
        IFAS & W1 & 0,14 & Akurasi packing diberi bobot tertinggi pada kelemahan karena muncul berulang pada observasi, pegawai, pelanggan, intern, dan pra-survei. Dampaknya langsung terhadap komplain, ulasan negatif, repeat order, dan kepercayaan pelanggan GrabFood. \\ \hline
        IFAS & W2 & 0,11 & Kecepatan layanan dan antrean lintas kanal diberi bobot tinggi karena berkaitan dengan waktu tunggu driver, kepuasan pelanggan, dan kapasitas saat promo atau jam ramai. Bobotnya sedikit di bawah W1 karena dampaknya paling kuat pada periode peak hour. \\ \hline
        IFAS & W3 & 0,09 & Update stok dan kesiapan menu penting karena dapat membuat pelanggan batal membeli atau pindah ke kompetitor. Bobotnya sedang karena bukti kuat, tetapi cakupannya lebih spesifik dibandingkan W1 dan W2. \\ \hline
        IFAS & W4 & 0,08 & Persepsi harga mahal dan ketergantungan promo penting karena memengaruhi keputusan checkout dan margin. Bobotnya sedang karena sebagian dampaknya dapat dikelola melalui bundling, komunikasi nilai porsi, dan pengaturan promo. \\ \hline
        IFAS & W5 & 0,07 & Informasi menu, label, segel, batas request, dan data pelanggan langsung diberi bobot lebih rendah karena lebih bersifat penguat pengalaman dan loyalitas jangka menengah. Faktor ini tetap penting, tetapi urgensinya berada setelah perbaikan akurasi, antrean, stok, dan harga. \\ \hline
        EFAS & O1 & 0,12 & Peluang permintaan makanan berat praktis dari pekerja, mahasiswa, dan pelanggan sekitar diberi bobot tinggi karena menjadi pasar inti Nasi Gerilya di GrabFood. Faktor ini langsung mendukung paket makan siang, makan malam, dan pelanggan yang mencari makanan mengenyangkan. \\ \hline
        EFAS & O2 & 0,12 & Promo, voucher, paket segmen, upselling, dan stamp digital diberi bobot tinggi karena berhubungan langsung dengan trial dan repeat order. Bobotnya sama dengan O1 karena peluang ini kuat, tetapi harus dikendalikan agar tidak menekan margin atau membebani operasional. \\ \hline
        EFAS & O3 & 0,10 & Tampilan digital, foto, rating, ulasan, deskripsi, dan menu digital penting untuk meningkatkan kepercayaan sebelum checkout. Bobotnya sedang-tinggi karena menjadi faktor konversi, tetapi berfungsi sebagai pendorong yang tetap membutuhkan produk dan proses yang stabil. \\ \hline
        EFAS & O4 & 0,09 & Masalah serupa pada kompetitor membuka ruang diferensiasi melalui ketelitian proses. Bobotnya sedang karena peluang ini bergantung pada kemampuan Nasi Gerilya membuktikan layanan yang lebih akurat dan cepat. \\ \hline
        EFAS & O5 & 0,08 & Media sosial, rekomendasi, ulasan pelanggan, dan database pelanggan memberi peluang awareness serta loyalitas di luar GrabFood. Bobotnya lebih rendah karena infrastruktur data dan program loyalitas masih perlu dibangun bertahap. \\ \hline
        EFAS & T1 & 0,13 & Reputasi digital kompetitor diberi bobot ancaman tertinggi karena pelanggan di platform mudah membandingkan rating, jumlah ulasan, foto, dan nama restoran sebelum membeli. Faktor ini sangat memengaruhi tahap pemilihan restoran. \\ \hline
        EFAS & T2 & 0,12 & Perbandingan harga, promo, ongkir, dan menu diberi bobot tinggi karena pelanggan GrabFood dapat berpindah cepat ketika harga akhir kompetitor lebih menarik. Bobotnya hampir sama dengan T1 karena memengaruhi keputusan checkout secara langsung. \\ \hline
        EFAS & T3 & 0,09 & Ongkir, biaya platform, kenaikan bahan baku, dan ketergantungan voucher menekan minat beli serta ruang margin. Bobotnya sedang karena dampaknya kuat, tetapi sebagian dapat dikendalikan melalui paket, batas promo, dan pengelolaan HPP. \\ \hline
        EFAS & T4 & 0,09 & Human error, makanan dingin, stok habis, antrean, dan item kurang dapat memicu ulasan negatif. Bobotnya sama dengan T3 karena berdampak pada reputasi, tetapi sumber masalahnya masih dapat dikurangi melalui perbaikan SOP internal. \\ \hline
        EFAS & T5 & 0,06 & Keunggulan spesifik kompetitor pada rasa, layanan cepat, sambal, porsi, dan menu diberi bobot terendah karena ancaman ini lebih spesifik dan sebagian dampaknya sudah tercermin pada reputasi digital serta perbandingan harga/promo. \\ \hline
    \end{longtable}
\endgroup

\begingroup
    \footnotesize
    \setlength{\tabcolsep}{2pt}
    \renewcommand{\arraystretch}{1.05}
    \begin{longtable}{|Z{0.85cm}|L{7.65cm}|C{1.05cm}|C{1.05cm}|C{1.05cm}|}
        \caption{Matriks IFAS Nasi Gerilya}
        \label{tab:ifas-bab4}\\
        \toprule
        \textbf{Kode} & \textbf{Faktor Internal} & \textbf{Bobot} & \textbf{Rating} & \textbf{Skor} \\ \hline
        \midrule
        \endfirsthead
        \multicolumn{5}{@{}l@{}}{\footnotesize\textit{Lanjutan Tabel \thetable}}\\
        \toprule
        \textbf{Kode} & \textbf{Faktor Internal} & \textbf{Bobot} & \textbf{Rating} & \textbf{Skor} \\ \hline
        \midrule
        \endhead
        \bottomrule
        \multicolumn{5}{r}{\footnotesize\textit{Bersambung ke halaman berikutnya}}\\
        \endfoot
        \bottomrule
        \multicolumn{5}{@{}L{\textwidth}@{}}{\footnotesize Sumber: Diolah peneliti berdasarkan hasil penelitian utama dan konteks pra-survei (2026).}\\
        \endlastfoot
        S1 & Rasa dan menu unggulan kuat. & 0,14 & 4 & 0,56 \\ \hline
        S2 & Porsi besar dan kemasan relatif rapi. & 0,10 & 4 & 0,40 \\ \hline
        S3 & Alur operasional online dan arah digitalisasi internal sudah memiliki checker, bill, nomor antrean, update lauk, serta rancangan sistem pendukung. & 0,11 & 3 & 0,33 \\ \hline
        S4 & Reputasi digital dan promo aktif cukup kuat. & 0,09 & 3 & 0,27 \\ \hline
        S5 & Respons keluhan sudah ada. & 0,07 & 3 & 0,21 \\ \hline
        W1 & Akurasi packing rawan saat ramai. & 0,14 & 1 & 0,14 \\ \hline
        W2 & Kecepatan layanan terhambat saat peak hour dan antrean lintas kanal. & 0,11 & 1 & 0,11 \\ \hline
        W3 & Update stok dan kesiapan menu belum sepenuhnya stabil. & 0,09 & 2 & 0,18 \\ \hline
        W4 & Harga di GrabFood dipersepsikan mahal dan bergantung pada promo. & 0,08 & 2 & 0,16 \\ \hline
        W5 & Informasi menu, segel, label, batas request, dan pengelolaan data pelanggan perlu dipertegas. & 0,07 & 2 & 0,14 \\ \hline
        \textbf{Total} & & \textbf{1,00} & & \textbf{2,50} \\ \hline
    \end{longtable}
\endgroup

\begingroup
    \footnotesize
    \setlength{\tabcolsep}{2pt}
    \renewcommand{\arraystretch}{1.05}
    \begin{longtable}{|Z{0.85cm}|L{7.65cm}|C{1.05cm}|C{1.05cm}|C{1.05cm}|}
        \caption{Matriks EFAS Nasi Gerilya}
        \label{tab:efas-bab4}\\
        \toprule
        \textbf{Kode} & \textbf{Faktor Eksternal} & \textbf{Bobot} & \textbf{Rating} & \textbf{Skor} \\ \hline
        \midrule
        \endfirsthead
        \multicolumn{5}{@{}l@{}}{\footnotesize\textit{Lanjutan Tabel \thetable}}\\
        \toprule
        \textbf{Kode} & \textbf{Faktor Eksternal} & \textbf{Bobot} & \textbf{Rating} & \textbf{Skor} \\ \hline
        \midrule
        \endhead
        \bottomrule
        \multicolumn{5}{r}{\footnotesize\textit{Bersambung ke halaman berikutnya}}\\
        \endfoot
        \bottomrule
        \multicolumn{5}{@{}L{\textwidth}@{}}{\footnotesize Sumber: Diolah peneliti berdasarkan hasil penelitian utama dan konteks pra-survei (2026).}\\
        \endlastfoot
        O1 & Segmen pekerja, mahasiswa, dan pelanggan sekitar membutuhkan makanan berat praktis. & 0,12 & 3 & 0,36 \\ \hline
        O2 & Promo, voucher, ongkir murah, paket segmen, dan stamp digital mendorong trial dan repeat order. & 0,12 & 3 & 0,36 \\ \hline
        O3 & Tampilan digital dan menu digital dapat diperkuat untuk menaikkan kepercayaan pelanggan baru. & 0,10 & 3 & 0,30 \\ \hline
        O4 & Kompetitor juga mengalami masalah stok, item tertinggal, dan driver menunggu. & 0,09 & 3 & 0,27 \\ \hline
        O5 & Media sosial, rekomendasi teman, ulasan, dan database pelanggan dapat memperluas awareness. & 0,08 & 3 & 0,24 \\ \hline
        T1 & Kompetitor memiliki reputasi digital kuat. & 0,13 & 2 & 0,26 \\ \hline
        T2 & Pelanggan mudah membandingkan harga, promo, ongkir, dan menu. & 0,12 & 2 & 0,24 \\ \hline
        T3 & Ongkir, biaya platform, kenaikan bahan baku, dan ketergantungan voucher menekan minat beli tanpa promo. & 0,09 & 2 & 0,18 \\ \hline
        T4 & Human error, antrean menumpuk, makanan dingin, stok habis, dan item kurang dapat memicu ulasan negatif. & 0,09 & 2 & 0,18 \\ \hline
        T5 & Kompetitor menonjolkan rasa, pelayanan cepat, sambal, porsi, dan menu spesifik. & 0,06 & 2 & 0,12 \\ \hline
        \textbf{Total} & & \textbf{1,00} & & \textbf{2,51} \\ \hline
    \end{longtable}
\endgroup

Total skor IFAS sebesar 2,50 menunjukkan posisi internal yang berada pada kondisi sedang: kekuatan produk, porsi, reputasi, SOP dasar, dan arah digitalisasi internal sudah ada, tetapi masih ditahan oleh kelemahan akurasi packing, antrean lintas kanal, kecepatan saat ramai, stok, dan persepsi harga. Total skor EFAS sebesar 2,51 menunjukkan respons eksternal yang juga sedang: peluang GrabFood, perilaku pelanggan, menu digital, dan loyalitas langsung cukup besar, tetapi tekanan kompetitor, biaya platform, kenaikan bahan baku, sensitivitas harga, dan risiko ulasan negatif perlu segera ditangani. Dalam pendekatan IFAS/EFAS, skor tersebut tidak dibaca sebagai angka statistik, tetapi sebagai alat bantu untuk menentukan arah strategi berdasarkan kepentingan relatif faktor internal dan eksternal \parencite{rangkuti2016analisisSWOT,david2023strategicmanagement}. Dengan demikian, strategi yang paling relevan bukan hanya memperbesar promosi, tetapi memperkuat proses layanan agar peluang digital tidak berubah menjadi keluhan.

\section{Matriks SWOT dan Alternatif Strategi}
Matriks SWOT disusun dengan mengombinasikan faktor kekuatan, kelemahan, peluang, dan ancaman yang sudah ditriangulasikan. SWOT digunakan sebagai alat untuk mencocokkan kemampuan internal dengan kondisi eksternal sehingga alternatif strategi tidak hanya bersifat deskriptif, tetapi mengarah pada tindakan yang dapat diprioritaskan \parencite{gurel2017swot,david2023strategicmanagement,rangkuti2016analisisSWOT}. Strategi yang dihasilkan disajikan pada Tabel~\ref{tab:swot-bab4}.

\begingroup
    \scriptsize
    \setlength{\tabcolsep}{2pt}
    \renewcommand{\arraystretch}{1.02}
    \begin{longtable}{|Z{2.2cm}|L{5.45cm}|L{5.45cm}|}
        \caption{Matriks SWOT Nasi Gerilya pada GrabFood}
        \label{tab:swot-bab4}\\
        \toprule
        \textbf{Faktor} & \textbf{Kekuatan (S)} & \textbf{Kelemahan (W)} \\ \hline
        \midrule
        \endfirsthead
        \multicolumn{3}{@{}l@{}}{\footnotesize\textit{Lanjutan Tabel \thetable}}\\
        \toprule
        \textbf{Faktor} & \textbf{Kekuatan (S)} & \textbf{Kelemahan (W)} \\ \hline
        \midrule
        \endhead
        \bottomrule
        \multicolumn{3}{r}{\footnotesize\textit{Bersambung ke halaman berikutnya}}\\
        \endfoot
        \bottomrule
        \multicolumn{3}{@{}L{\textwidth}@{}}{\footnotesize Sumber: Diolah peneliti berdasarkan hasil IFAS dan EFAS (2026).}\\
        \endlastfoot
        \textbf{Peluang (O)} & \textbf{SO:} Mengangkat menu hero dendeng, ayam pop, dan ayam goreng ke paket GrabFood untuk pekerja, mahasiswa, dan pelanggan sekitar; memperkuat foto, deskripsi, rating, promo, menu digital, dan stamp digital agar pelanggan baru maupun pelanggan ulang lebih cepat percaya. & \textbf{WO:} Menggunakan promo dan paket segmen hanya pada menu yang stoknya siap; memperjelas isi paket, kuah, sambal, segel, label, serta menerapkan sistem antrean dan checklist packing untuk mengurangi item kurang. \\ \hline
        \textbf{Ancaman (T)} & \textbf{ST:} Mempertahankan diferensiasi rasa, porsi, respons keluhan, dan kemauan berbenah melalui sistem digital untuk menghadapi kompetitor dengan reputasi tinggi; menonjolkan menu yang punya karakter lebih jelas dibanding pesaing. & \textbf{WT:} Memperbaiki SOP peak hour, sistem antrean, update stok, estimasi tunggu driver, validasi nomor pesanan, dan evaluasi biaya promo agar kesalahan operasional serta tekanan margin tidak memicu ulasan negatif atau perpindahan pelanggan ke kompetitor. \\ \hline
    \end{longtable}
\endgroup

\section{Prioritas Strategi Pemasaran untuk Meningkatkan Penjualan}
Berdasarkan matriks SWOT, prioritas strategi disusun dari strategi yang memiliki hubungan paling langsung dengan penjualan, pembelian ulang, dan daya saing di GrabFood. Prioritas tidak hanya mengikuti jumlah faktor yang terkait, tetapi juga mempertimbangkan logika pemasaran jasa bahwa kualitas proses dan pengalaman pelanggan harus mendukung promosi agar pembelian pertama dapat berubah menjadi pembelian ulang \parencite{wirtz2024servicesmarketing,kotler2023principlesmarketing}. Prioritas strategi disajikan pada Tabel~\ref{tab:prioritas-strategi-bab4}.

\begingroup
    \scriptsize
    \setlength{\tabcolsep}{2pt}
    \renewcommand{\arraystretch}{1.02}
    \begin{longtable}{|Z{0.9cm}|L{3.9cm}|L{3.25cm}|L{2.65cm}|L{2.25cm}|}
        \caption{Prioritas Strategi Pemasaran Nasi Gerilya}
        \label{tab:prioritas-strategi-bab4}\\
        \toprule
        \textbf{No.} & \textbf{Strategi} & \textbf{Dasar Data} & \textbf{Dampak yang Diharapkan} & \textbf{Alasan Kelayakan} \\ \hline
        \midrule
        \endfirsthead
        \multicolumn{5}{@{}l@{}}{\footnotesize\textit{Lanjutan Tabel \thetable}}\\
        \toprule
        \textbf{No.} & \textbf{Strategi} & \textbf{Dasar Data} & \textbf{Dampak yang Diharapkan} & \textbf{Alasan Kelayakan} \\ \hline
        \midrule
        \endhead
        \bottomrule
        \multicolumn{5}{r}{\footnotesize\textit{Bersambung ke halaman berikutnya}}\\
        \endfoot
        \bottomrule
        \multicolumn{5}{@{}L{\textwidth}@{}}{\footnotesize Sumber: Diolah peneliti berdasarkan hasil penelitian utama (2026).}\\
        \endlastfoot
        1 & Membuat sistem antrean lintas kanal, checklist packing GrabFood, dan SOP peak hour untuk dine in, ojek online, take away, item terpisah, kuah, sambal, catatan khusus, nomor pesanan, serta estimasi tunggu driver. & W1, W2, W3, T4; TW-AD-05, TS-PL-09, TO-GF-08, W-BC-05, W-BC-15, W-BC-20 & Mengurangi item kurang, komplain, driver menunggu, antrean menumpuk, dan ulasan negatif. & Masalah antrean terlihat langsung di lapangan dan sistem antrean sudah menjadi rekomendasi kuat dari pendamping internal. \\ \hline
        2 & Mengangkat paket menu hero berbasis segmen: paket dendeng, paket ayam pop kantor, paket ayam goreng anak kos, dan bundling minum/sambal/kuah. & S1, S2, O1, O2; TS-PL-03, TS-PL-04, TS-PL-06, W-BC-07, W-BC-09, PS-WPM-011, PS-GF-031--PS-GF-037, PS-PROD-003, PS-PROD-014 & Meningkatkan trial, repeat order, AOV, dan persepsi value. & Sesuai dengan menu yang sudah disukai pelanggan dan dapat dipromosikan melalui GrabFood. \\ \hline
        3 & Memperbaiki tampilan digital dan menu digital: foto menu hero, detail isi paket, informasi kuah/sambal, ketersediaan stok, label/segel, posisi menu unggulan, dan pengalaman pemesanan mandiri. & S4, W5, O3, T2; TS-PL-07, TS-PL-08, W-BC-05, W-BC-11 & Meningkatkan kepercayaan pelanggan baru, menurunkan keraguan sebelum checkout, dan mengurangi beban pencatatan manual. & Perubahan dapat dilakukan bertahap melalui pembaruan konten toko digital dan sistem menu digital. \\ \hline
        4 & Menyinkronkan promo dengan kesiapan stok, kapasitas dapur, HPP, dan margin, terutama sebelum jam makan siang, malam, akhir pekan, hari Minggu, dan hari raya. & W2, W3, O2, T3; TW-AD-10, TW-KR-03, W-BC-10, W-BC-17, PS-WPM-014, PS-WPM-015, PS-WPM-017, PS-WPM-018 & Meningkatkan volume order tanpa memperbesar risiko keterlambatan, stok habis, atau tekanan margin. & Dapat dilakukan melalui briefing harian, pemantauan menu promo, dan pembatasan diskon berdasarkan kesiapan menu. \\ \hline
        5 & Mengaktifkan stamp digital dan database pelanggan untuk program loyalitas dine in/take away serta komunikasi ulang kepada pelanggan. & W5, O2, O5, T3; W-BC-05, W-BC-12, W-BC-16, W-BC-18 & Meningkatkan repeat order dan mengurangi ketergantungan penuh pada voucher platform. & Relevan karena pelanggan offline belum memiliki promo loyalitas sendiri dan data pelanggan dapat dipakai untuk evaluasi. \\ \hline
    \end{longtable}
\endgroup

\section{Pembahasan}

\subsection{Pembahasan Strategi Pemasaran Berdasarkan STP dan 7P}
Strategi pemasaran Nasi Gerilya pada GrabFood sudah memiliki dasar yang kuat pada produk, porsi, reputasi digital, dan penggunaan promo. Pelanggan mengenali Nasi Gerilya melalui pencarian aplikasi, rating, foto menu, dan menu yang terlihat menarik. Hal ini sesuai dengan konsep pemasaran digital bahwa keputusan pelanggan dalam platform tidak hanya dipengaruhi produk, tetapi juga informasi digital yang dapat dibandingkan sebelum pembelian \parencite{kotler2023principlesmarketing,chaffey2022digitalmarketing,sophia2025rendanggadih}.

Namun, 7P menunjukkan bahwa unsur \textit{process} menjadi titik penentu. Produk yang enak dan porsi yang besar dapat kehilangan dampaknya ketika pesanan kurang lengkap, driver menunggu terlalu lama, stok tidak diperbarui, antrean dine in/ojek online/take away menumpuk, atau harga akhir terasa terlalu mahal. Dalam pemasaran jasa, proses, orang, dan bukti fisik/digital menjadi bagian dari kualitas layanan yang menentukan kepuasan, kepercayaan, dan kemungkinan pembelian ulang \parencite{wirtz2024servicesmarketing,kotler2023principlesmarketing}. Dengan kata lain, strategi pemasaran tidak cukup dilakukan lewat promo dan foto menu, tetapi harus ditopang oleh SOP pesanan, sistem antrean, dan data pelanggan yang konsisten.

\subsection{Pembahasan Faktor Internal dan Eksternal}
Kekuatan utama Nasi Gerilya berada pada rasa, menu unggulan, porsi besar, reputasi digital, dan kemauan untuk berbenah melalui sistem digital. Kelemahan utama berada pada akurasi packing, kecepatan layanan saat ramai, antrean lintas kanal, stok, pengelolaan data pelanggan, dan persepsi harga. Faktor eksternal menunjukkan peluang yang besar dari kebiasaan pelanggan memesan makanan berat melalui GrabFood, terutama saat sibuk atau malas keluar, serta peluang membangun repeat order melalui stamp digital. Akan tetapi, ancaman dari kompetitor juga nyata karena pelanggan membandingkan rating, harga, ongkir, promo, foto, porsi, dan menu dengan restoran lain. Pola ini sejalan dengan kerangka SWOT yang menempatkan faktor internal dan eksternal sebagai dasar pencocokan strategi, bukan sekadar daftar temuan lapangan \parencite{david2023strategicmanagement,gurel2017swot,rangkuti2016analisisSWOT}.

Temuan kompetitor memperkuat bahwa isu item kurang, stok, driver menunggu, dan konsistensi porsi bukan masalah Nasi Gerilya saja. Justru karena masalah tersebut bersifat umum, Nasi Gerilya dapat menjadikannya ruang diferensiasi. Diferensiasi yang efektif tidak hanya berasal dari keunikan produk, tetapi juga dari kemampuan memberikan nilai dan pengalaman yang lebih dapat dipercaya dibandingkan alternatif lain \parencite{kotler2023principlesmarketing,wirtz2024servicesmarketing}. Jika Nasi Gerilya mampu membuktikan pesanan lebih akurat dan stok lebih jelas, maka kekuatan rasa dan porsi akan lebih mudah berubah menjadi pembelian ulang.

\subsection{Pembahasan Strategi Prioritas}
Strategi prioritas pertama adalah memperkuat proses, bukan menambah promo. Hal ini karena promo memang mampu menaikkan pesanan, tetapi juga dapat memperbesar beban operasional. Jika promo tidak disertai kesiapan stok, packing, dan sistem antrean, maka peningkatan order dapat berubah menjadi keterlambatan, item kurang, antrean menumpuk, atau ulasan negatif. Oleh karena itu, sistem antrean, checklist packing, dan SOP peak hour menjadi strategi dasar. Urutan ini sejalan dengan pemasaran jasa dan pemasaran digital karena promosi yang efektif perlu ditopang oleh pengalaman layanan yang konsisten agar tidak merusak rating, ulasan, dan kepercayaan pelanggan \parencite{wirtz2024servicesmarketing,chaffey2022digitalmarketing}.

Strategi kedua adalah mengemas menu unggulan dalam paket yang sesuai segmen. Data pelanggan memperlihatkan bahwa tiap segmen memiliki alasan berbeda: pekerja membutuhkan makan siang cepat, mahasiswa membutuhkan porsi besar dan promo, pelanggan sekitar memperhatikan ongkir, sedangkan pelanggan yang mengejar rasa tertarik pada dendeng, ayam pop, dan ayam goreng. Wawancara intern \textit{business consultant} juga menegaskan bahwa ayam pop layak dijadikan hero menu. Paket yang jelas dapat membantu pelanggan memilih lebih cepat dan membuat harga terasa lebih wajar. Secara teoritis, strategi ini mempertemukan targeting, positioning, dan penciptaan nilai pelanggan melalui penawaran yang lebih mudah dipahami oleh segmen sasaran \parencite{kotler2023principlesmarketing,sophia2025rendanggadih}.

Strategi lanjutan adalah menghubungkan pemasaran platform dengan data pelanggan sendiri. Stamp digital dapat membantu Nasi Gerilya membangun hubungan langsung dengan pelanggan dine in dan take away, sementara menu digital dapat mengurangi beban pencatatan manual. Dengan demikian, strategi pemasaran tidak hanya bergantung pada voucher GrabFood, tetapi juga pada loyalitas pelanggan yang dapat diukur melalui repeat order, jumlah antrean tertangani, rating, komplain, dan efektivitas sistem. Pengelolaan hubungan pelanggan dan data digital penting karena pelanggan yang sudah pernah membeli dapat dikelola kembali melalui komunikasi, insentif, dan evaluasi perilaku pembelian \parencite{chaffey2022digitalmarketing,kotler2021marketingmanagement}.

\subsection{Implikasi Manajerial}
Secara manajerial, Nasi Gerilya perlu mengelola GrabFood sebagai sistem layanan, bukan hanya kanal transaksi. Implikasi ini konsisten dengan pandangan bahwa strategi pemasaran perlu diterjemahkan ke proses operasional, pengendalian kualitas layanan, dan evaluasi kinerja yang dapat ditindaklanjuti \parencite{wirtz2024servicesmarketing,david2023strategicmanagement}. Pertama, setiap pesanan perlu melalui checklist sederhana: nomor pesanan, lauk utama, item terpisah, kuah, sambal, add-on, catatan khusus, alat makan, segel, dan nama/label pesanan. Kedua, sistem antrean perlu memisahkan alur dine in, ojek online, dan take away agar beban jam ramai lebih terbaca. Ketiga, menu promo harus diinformasikan ke dapur sebelum aktif agar stok siap dan margin tetap terjaga. Keempat, halaman GrabFood serta menu digital perlu menonjolkan menu hero dan memperjelas isi paket. Kelima, stamp digital dan database pelanggan perlu dipakai untuk membangun repeat order di luar ketergantungan pada voucher platform. Keenam, pengelola perlu memantau jumlah pesanan, omzet, waktu pelayanan, antrean tertangani, repeat order, rating, item kurang, menu habis, komplain, dan efektivitas sistem sebagai data rutin untuk memperbaiki layanan.
